\section{深层思考：严谨性的迁移}

\subsection{超越 Markdown 规则文件}

Birgitta 指出，OpenAI 团队描述的工作远比``生成和维护一堆 Markdown 规则文件''要复杂得多。他们为 harness 的确定性部分构建了\textbf{大量的工具}，他们的上下文工程不仅涉及知识库的整理，还涉及大量的\textbf{设计工作}。

这一观察击中了当前 AI 编程工具使用中的一个常见误区：许多团队认为编写 `.cursorrules` 或类似的 Markdown 规则文件就足够了。但 OpenAI 的实践表明，真正有效的 harness 需要工程投入（engineering effort），而非仅仅是提示工程（prompt engineering）。

\subsection{``严谨性的迁移''（Relocating Rigor）}

OpenAI 团队表示：``我们目前最困难的挑战集中在设计环境、反馈循环和控制系统上。'' 这让 Birgitta 联想到了 Chad Fowler 最近关于 ``Relocating Rigor''（严谨性的迁移）的文章。

\begin{importantbox}{严谨性的迁移}
在 AI 辅助编程时代，工程严谨性并没有消失，而是从一个层面迁移到了另一个层面：
\begin{itemize}[leftmargin=1.5em]
    \item \textbf{之前}：严谨性体现在手动编写代码------命名规范、设计模式、代码审查
    \item \textbf{之后}：严谨性迁移到设计环境、反馈循环和控制系统------即 harness 的构建和维护
\end{itemize}
这不是严谨性的减少，而是严谨性的重新定位。听到关于``严谨性应该去向何处''的具体想法和经验，比仅仅寄希望于``更好的模型''会神奇地解决可维护性问题要令人振奋得多。
\end{importantbox}

\subsection{本章小结}

有效的 harness 需要真正的工程投入，而非仅靠提示工程。工程严谨性从手动编码层面迁移到了环境设计、反馈循环和控制系统的构建上。这种``严谨性的迁移''是 AI 编程时代最重要的认知转变之一。

\newpage

